\documentclass[12pt,a4paper,openany,oneside]{book}

% --- Paquetes ---
\usepackage[utf8]{inputenc}
\usepackage[spanish, es-noshorthands]{babel}
\usepackage{amsmath, amsfonts, amssymb}
\usepackage{graphicx}
\usepackage{geometry}
\usepackage{hyperref}
\usepackage{booktabs}
\usepackage{fancyhdr}
\usepackage{listings}
\usepackage{xcolor}
\usepackage{tikz}
\usepackage{pgfplots}
\usepackage{colortbl}
\usepackage{multirow}
\usepackage{hhline}
\usetikzlibrary{arrows.meta, positioning, shapes.geometric, calc, shapes, chains, scopes, shadows}
\pgfplotsset{compat=1.18}

\geometry{margin=2.5cm}

% --- Colores y estilos para código Python ---
\definecolor{codegreen}{rgb}{0,0.6,0}
\definecolor{codegray}{rgb}{0.5,0.5,0.5}
\definecolor{codepurple}{rgb}{0.58,0,0.82}
\definecolor{backcolour}{rgb}{0.95,0.95,0.92}

\lstset{
    language=Python,
    backgroundcolor=\color{backcolour},   
    commentstyle=\color{codegreen},
    keywordstyle=\color{blue},
    numberstyle=\tiny\color{codegray},
    stringstyle=\color{codepurple},
    basicstyle=\ttfamily\small,
    breaklines=true,
    frame=single,
    showstringspaces=false
}

% --- Estilo de cabecera y pie ---
\pagestyle{fancy}
\fancyhf{}
\renewcommand{\headrulewidth}{0.4pt}
\renewcommand{\footrulewidth}{0.4pt}
\fancyhead[L]{\small Carlos César Sánchez Coronel}
\fancyhead[R]{\small Sesión 13: Prompt Engineering y Optimización}
\fancyfoot[L]{\small Carlos César Sánchez Coronel}
\fancyfoot[C]{\thepage}

% --- Portada ---
\title{
    \textbf{Sesión 13} \\ 
    \Huge \textbf{PROMPT ENGINEERING Y OPTIMIZACIÓN DE MODELOS} \\
    \Large In-Context Learning, LoRA y Cuantización
}
\author{
    \small Por \\ \vspace{0.2cm}
    \Large \textsc{Carlos César Sánchez Coronel}
}
\date{2026}

\begin{document}

\maketitle
\tableofcontents
\addcontentsline{toc}{chapter}{Índice general}

% ------------------------------------------------------------------
% CAPÍTULO 13: PROMPT ENGINEERING Y OPTIMIZACIÓN
% ------------------------------------------------------------------
\chapter{Prompt Engineering y Optimización de Modelos}

La capacidad de un modelo de lenguaje no solo depende de sus parámetros, sino de cómo interactuamos con él y cómo optimizamos sus recursos para entornos de producción. Esta sesión aborda el arte de diseñar entradas efectivas y las técnicas de ingeniería avanzadas para reducir el costo computacional y mejorar la especialización del modelo.

\section{Logro de la sesión}
Dominar técnicas avanzadas de \textit{prompt engineering} y aplicar métodos de optimización de modelos como LoRA (Low-Rank Adaptation) y cuantización para lograr una inferencia eficiente y un ajuste fino con recursos limitados.

\section{Prompt Engineering y In-Context Learning}

El \textbf{In-context Learning (ICL)} es la capacidad de los LLMs para aprender a resolver tareas simplemente observando ejemplos en el prompt, sin cambiar sus pesos internos.

\subsection{Estrategias Principales}
\begin{itemize}
    \item \textbf{Zero-shot}: El modelo recibe una instrucción sin ejemplos previos. Depende totalmente del conocimiento pre-entrenado.
    \item \textbf{Few-shot}: Se proporcionan algunos pares de entrada-salida (ejemplos) para guiar al modelo en el formato y estilo deseado.
    \item \textbf{Chain-of-Thought (CoT)}: Se le pide al modelo que "piense paso a paso". Esto descompone problemas complejos en razonamientos intermedios, mejorando drásticamente el desempeño en tareas lógicas o matemáticas.
\end{itemize}

\textbf{Ejemplo de CoT}:
\begin{quote}
    \textit{"P: Si tengo 3 manzanas y compro 2 docenas más, ¿cuántas tengo? Pensémoslo paso a paso. \\
    R: Primero, una docena son 12. Dos docenas son 24. 24 + 3 iniciales = 27. Total: 27."}
\end{quote}

\section{Fine-Tuning Eficiente: LoRA (Low-Rank Adaptation)}

El ajuste fino tradicional requiere actualizar miles de millones de parámetros, lo cual es prohibitivo. LoRA propone que los cambios en los pesos durante la adaptación tienen un "rango intrínseco bajo".

\subsection{Arquitectura de LoRA}
En lugar de modificar la matriz original de pesos $W_0 \in \mathbb{R}^{d \times k}$, LoRA entrena dos matrices pequeñas $A$ y $B$ tales que:
\[ W = W_0 + \Delta W = W_0 + B \cdot A \]
donde $A \in \mathbb{R}^{r \times k}$ y $B \in \mathbb{R}^{d \times r}$, siendo el rango $r \ll \min(d, k)$.

\begin{center}
\begin{tikzpicture}[
    node distance=1cm,
    matrix_block/.style={rectangle, draw, fill=blue!10, minimum width=2cm, minimum height=1.5cm, align=center, font=\small},
    small_matrix/.style={rectangle, draw, fill=orange!20, align=center, font=\scriptsize},
    arrow/.style={-Stealth, thick}
]
    % Camino Original
    \node (input) {Entrada $x$};
    \node[matrix_block, right=of input] (W0) {$W_0$ \\ (Congelado)};
    
    % Camino LoRA
    \node[small_matrix, above right=0.5cm and 0.5cm of input, minimum width=1.5cm, minimum height=0.4cm] (A) {Matriz $A$};
    \node[small_matrix, right=0.5cm of A, minimum width=0.4cm, minimum height=1.5cm] (B) {Matriz $B$};
    
    % Suma y Salida
    \node[circle, draw, right=3.5cm of W0] (sum) {+};
    \node[right=of sum] (output) {Salida $h$};

    % Conexiones
    \draw[arrow] (input) -- (W0);
    \draw[arrow] (W0) -- (sum);
    \draw[arrow] (input.east) -- ++(0.2,0) |- (A.west);
    \draw[arrow] (A) -- (B);
    \draw[arrow] (B.east) -| (sum.north);
    \draw[arrow] (sum) -- (output);

\end{tikzpicture}
\captionof{figure}{Mecanismo de LoRA: Solo las matrices $A$ y $B$ se entrenan.}
\end{center}

\section{Cuantización: Inferencia a Bajo Costo}

La cuantización reduce la precisión de los pesos (ej. de FP32 a INT8 o FP16) para ahorrar memoria y acelerar el procesamiento.

\subsection{Tipos de Precisiones}
\begin{table}[h]
\centering
\begin{tabular}{l l l}
\toprule
\textbf{Formato} & \textbf{Bits} & \textbf{Impacto} \\ \midrule
\textbf{FP32} & 32 bits & Precisión completa, alto consumo de VRAM. \\
\textbf{FP16 / BF16} & 16 bits & Estándar para entrenamiento moderno. \\
\textbf{INT8} & 8 bits & Reduce el tamaño del modelo al 25\% del original. \\
\textbf{4-bit (NF4)} & 4 bits & Utilizado en QLoRA para modelos masivos en GPUs domésticas. \\ \bottomrule
\end{tabular}
\caption{Comparativa de formatos de cuantización.}
\end{table}

\section{Exportación y Despliegue: ONNX}

Para llevar modelos a producción fuera de Python (ej. C++, navegadores o móviles), se utiliza \textbf{ONNX} (Open Neural Network Exchange). Permite la interoperabilidad entre frameworks como PyTorch y TensorFlow.

\section{Preparación para Entrevistas}

\subsection{Pregunta Avanzada: "¿Por qué LoRA es más eficiente que el fine-tuning completo?"}
\textbf{Respuesta}: LoRA reduce el número de parámetros entrenables en un 10,000\% aproximadamente. Al mantener los pesos originales congelados, no hay "olvido catastrófico" y se pueden intercambiar adaptadores (matrices $A$ y $B$) para diferentes tareas sin recargar el modelo base.

\subsection{El Dilema de la Cuantización}
Es fundamental saber que la cuantización reduce el tamaño y acelera la inferencia, pero a costa de una pequeña pérdida de precisión (perplejidad). En modelos grandes (>13B), esta pérdida es casi imperceptible.

\section{Implementación Conceptual (Python)}

\begin{lstlisting}[language=Python]
from peft import LoraConfig, get_peft_model
from transformers import AutoModelForCausalLM

# 1. Cargar modelo base
model = AutoModelForCausalLM.from_pretrained("base-model")

# 2. Configurar LoRA
config = LoraConfig(
    r=8, # Rango
    lora_alpha=32,
    target_modules=["q_proj", "v_proj"],
    lora_dropout=0.05,
    bias="none"
)

# 3. Preparar modelo para entrenamiento eficiente
lora_model = get_peft_model(model, config)
lora_model.print_trainable_parameters()
\end{lstlisting}

\section{Resumen de la Sesión}
\begin{itemize}
    \item \textbf{Prompt Engineering}: Uso de Zero/Few-shot y Chain-of-Thought para maximizar el rendimiento sin reentrenar.
    \item \textbf{LoRA}: Adaptación de bajo rango para entrenar modelos gigantes con poca memoria.
    \item \textbf{Cuantización}: Técnica esencial para comprimir modelos y habilitar inferencia rápida en hardware limitado.
    \item \textbf{ONNX}: Estándar industrial para portabilidad de modelos.
\end{itemize}

\end{document}